\section{Orientations: The Religious Question}

\begin{quotex}
The \emph{nature} of man is to be a cognitive, religious and sociable animal. All experience teaches us this; and, to my knowledge, nothing has contradicted this experience.
\flright{\textsc{Joseph de Maistre}, \emph{Study of Sovereignty}}
\end{quotex}

Evola regarded Joseph de Maistre as standing on the same side of the barricade as himself. And, like Maistre, he does not deny the obvious --- man is a cognitive, religious and sociable animal. It is only the characteristics of knowledge, religion, and society that is in question.

In 1950, Evola wrote a short pamphlet, \emph{Orientations}, consisting of eleven fundamental points to guide his young followers who were out to form a new right movement after the collapse of World War II. The eleventh and final point deals with the religious question.

\begin{quotex}
A religious factor is necessary as the background for a true heroic conception of life which must be essential for our battle deployment. It is necessary to feel in oneself the evidence that there is a higher life beyond this terrestrial life because only those who feel that way possess an intangible and inner directed strength, only they will be capable of an absolute impulse --- while, when this is lacking, defying death and taking no account of one's own life is possible only in sporadic moments of exaltation or in an outburst of irrational forces: nor is there discipline that can be justified, in the individual, with a higher and self-sufficient significance.
\end{quotex}

\flrightit{Posted on 2008-08-25 by Cologero}

\begin{center}* * *\end{center}

\begin{footnotesize}
\begin{sffamily}
\texttt{GF on 2010-05-25 at 19:48 said:}

In the Hindu Doctrines, chapter four, Guenon's definition of religion pointedly includes sociability and cognitivity. It is also worth noting that etymologically the word religion has two meanings: that which binds men to men and that which binds men to a higher principle. Clearly the social bond, in its deepest sense, is only possible by means of a shared principle. What's more Guenon says that metaphysical incapacity ``brings a strange confusion of thought as its fatal and unquestionable consequence.'' (58) So it seems that mans cognitive and sociable natures are wrapped up in his religion (in the principled sense).

While on ch. 4, I have to ask: who's evaluation of Rome is more correct, Evola's or Guenon's?

\hfill

\texttt{Cologero on 2010-05-26 at 23:27 said:}

Do you have something specific about Rome?

Evola would regard Rome from the perspective of Action. Also, the Emperor being both spiritual as well as political leader. He considered the Roman mystery cults to be authentic sources of initiation.

Guenon's perspective is contemplation, and the Roman religion was not as anvanced as the Vedanta. Guenon thought the political authority should be under the spiritual authority.

\hfill

\texttt{GFJF on 2010-05-29 at 02:20 said:}

That answers my question, actually. Should have seen it for myself.

\hfill

\end{sffamily}
\end{footnotesize}

